\section*{第 \ref{ch:g11:titration} 章——滴定}

\begin{solution}{exo:g11:titration:1}
\omterm{def:g11:titration:titration}{滴定剂}是浓度准确已知的溶液，装于
\omterm{def:g11:titration:titration}{滴定}管中；\omterm{def:g11:titration:titration}{被滴定溶液}位于锥形瓶中，含有
待测物质。\omterm{def:g11:titration:equivalence}{滴定终点}时，所加\omterm{def:g11:titration:titration}{滴定剂}与待测
物质恰按方程式中的比例反应，二者均被耗尽。
\end{solution}

\begin{solution}{exo:g11:titration:2}
$n = 0.0200 \times 0.0125 = \qty{2.50e-4}{mol}$.
\end{solution}

\begin{solution}{exo:g11:titration:3}
$V_E = \qty{14.3}{mL}$；\omterm{def:g10:reaction-progress-table:extent}{进度表}中亚铁\omterm{def:g9:ions:ion}{离子}一行的起始量
为 \qty{14.3e-4}{mol}，即 \qty{1.43e-3}{mol}。
\end{solution}

\begin{solution}{exo:g11:titration:4}
\omterm{def:g11:titration:titration}{滴定剂}本身有颜色（紫色），其生成物
几乎无色：一旦高锰酸根过量，颜色便
出现在锥形瓶中。
\end{solution}

\begin{solution}{exo:g11:titration:5}
5：方程式中每个高锰酸根\omterm{def:g9:ions:ion}{离子}消耗 5 个亚铁\omterm{def:g9:ions:ion}{离子}，\omterm{def:g11:titration:equivalence}{滴定终点}时
\omterm{def:g7:chemical-reactions:reactant}{反应物}已按此比例恰好\omterm{def:g2:mixing-and-dissolving:mix}{混合}。
%
\end{solution}

\begin{solution}{exo:g11:titration:6}
$n(\ce{MnO4^-}) = 0.0200 \times 0.00840 = \qty{1.68e-4}{mol}$;
$n(\ce{Fe^{2+}}) = 5 \times \num{1.68e-4} = \qty{8.40e-4}{mol}$;
$c = \num{8.40e-4} / 0.0100 = \qty{0.0840}{mol/L}$;
$C_m = 0.0840 \times 55.8 = \qty{4.69}{g/L}$.
\end{solution}

\begin{solution}{exo:g11:titration:7}
\omterm{def:g10:the-mole:amount}{物质的量}之比为 1:1：$n = 0.0100 \times 0.0114 = \qty{1.14e-4}{mol}$，在
\qty{10.0}{mL} 中，因此 \qty{100.0}{mL} 中含 \qty{1.14e-3}{mol}，即
$\num{1.14e-3} \times 176.0 = \qty{0.201}{g}$：约 \qty{200}{mg}。
\end{solution}

\begin{solution}{exo:g11:titration:8}
每滴\omterm{def:g11:titration:titration}{滴定剂}不会立即反应，\omterm{def:g11:titration:titration}{滴定剂}的颜色会
在\omterm{def:g11:titration:equivalence}{滴定终点}之前持续存在，终点会被看得过早（或
需要每滴之后等待数分钟）。
\end{solution}

\begin{solution}{exo:g11:titration:9}
偏大：她已超过\omterm{def:g11:titration:equivalence}{滴定终点}，加入的高锰酸根过量。
接近\omterm{def:g11:titration:equivalence}{滴定终点}时，她本应逐滴加入\omterm{def:g11:titration:titration}{滴定剂}，
边摇动边加，并在出现持久的淡粉色时停止。
\end{solution}

\begin{solution}{exo:g11:titration:10}
$n(\ce{Fe^{2+}}) = 5 \times 0.0200 \times 0.0130 = \qty{1.30e-3}{mol}$ 位于
\qty{10.0}{mL} 溶液中：\omterm{def:g10:concentration-and-dilution:dilution}{稀释}后溶液的浓度 $c = \qty{0.130}{mol/L}$，而
原溶液的浓度为 $10 \times 0.130 = \qty{1.30}{mol/L}$。
\end{solution}

\begin{solution}{exo:g11:titration:11}
$0.0630 \times 55.8 = \qty{3.52}{g}$.
\end{solution}

\begin{solution}{exo:g11:titration:12}
$0.1 / 7.2 \approx \qty{1.4}{\percent}$，而
取 \qty{14.3}{mL} 时仅为 \qty{0.7}{\percent}。应
取更大体积（或更浓的样品）的待测溶液进行\omterm{def:g11:titration:titration}{滴定}，或使用更稀的\omterm{def:g11:titration:titration}{滴定剂}。
\end{solution}

\begin{solution}{exo:g11:titration:13}
$n(\ce{MnO4^-}) = 0.0200 \times 0.0176 = \qty{3.52e-4}{mol}$;
$n(\ce{H2O2}) = \frac{5}{2} \times \num{3.52e-4} = \qty{8.80e-4}{mol}$ 在
\qty{10.0}{mL} 中：\omterm{def:g10:concentration-and-dilution:dilution}{稀释}后溶液中为 \qty{0.0880}{mol/L}，
消毒液中为 \qty{0.880}{mol/L}，即
$0.880 \times 34.0 = \qty{29.9}{g/L}$（约为 3 \% 的溶液）。
\end{solution}

\begin{solution}{exo:g11:titration:14}
$n(\ce{MnO4^-}) = \num{1.4e-3} / 5 = \qty{2.8e-4}{mol}$，需在
约 \qty{0.015}{L} 溶液中配成：$c \approx \qty{0.019}{mol/L}$，故取
\qty{0.0200}{mol/L} 的溶液。若浓度提高十倍，$V_E$ 将约为
\qty{1.4}{mL}，\omterm{def:g11:titration:titration}{滴定}管误差将达到约
\qty{7}{\percent}。
\end{solution}

\begin{solution}{exo:g11:titration:15}
所需氢\omterm{def:g9:ions:ion}{离子}的\omterm{def:g10:the-mole:amount}{物质的量}为 $8 \times 0.0200 \times 0.0143 = \qty{2.29e-3}{mol}$。
\qty{10}{mL} 浓度为 \qty{1.0}{mol/L} 的酸中含 \qty{1.0e-2}{mol}，足够，超出
所需四倍以上。若无酸，本章中的反应无法
按所写方式进行（氢\omterm{def:g9:ions:ion}{离子}是\omterm{def:g7:chemical-reactions:reactant}{反应物}），而高锰酸根
将以不同方式反应。
\end{solution}

\begin{solution}{pb:g11:titration:1}
\textbf{1.} \ce{MnO4^-}/\ce{Mn^{2+}} 与 \ce{Fe^{3+}}/\ce{Fe^{2+}} 是\omterm{def:g11:redox:couple}{氧化还原电对}。
其中\omterm{def:g11:redox:oxidant}{氧化剂}是高锰酸根\omterm{def:g9:ions:ion}{离子}，\omterm{def:g11:redox:oxidant}{还原剂}是亚铁\omterm{def:g9:ions:ion}{离子}。

\textbf{2.} \ce{MnO4^- + 8H+ + 5e- -> Mn^{2+} + 4H2O};
\ce{Fe^{2+} -> Fe^{3+} + e-}.

\textbf{3.} 将第二式乘以 5：
\ce{MnO4^- + 8H+ + 5Fe^{2+} -> Mn^{2+} + 5Fe^{3+} + 4H2O}。左侧电荷
为 $-1 + 8 + 10 = +17$；右侧为 $+2 + 15 = +17$。

\textbf{4.} 氢\omterm{def:g9:ions:ion}{离子}是\omterm{def:g7:chemical-reactions:reactant}{反应物}：该反应需要
\omterm{def:g9:acids-bases-ph:acidic}{酸性溶液}。

\textbf{5.} 该反应完全、迅速，且是高锰酸根
在锥形瓶中唯一发生的反应；紫色的高锰酸根在\omterm{def:g11:titration:equivalence}{滴定终点}前已耗尽，
\omterm{def:g11:titration:equivalence}{滴定终点}后则使锥形瓶呈粉色。

\textbf{6.} 在\omterm{def:g11:titration:titration}{滴定}管中；视线与弯月面底部齐平时读数，
前后各读一次，精确到 \qty{0.05}{mL}。

\textbf{7.} 淡黄色（三价铁\omterm{def:g9:ions:ion}{离子}）；每滴紫色的高锰酸根
都会立即被仍存在的亚铁\omterm{def:g9:ions:ion}{离子}耗尽。

\textbf{8.} $0.0200 \times 0.0143 = \qty{2.86e-4}{mol}$.

\textbf{9.} $n(\ce{Fe^{2+}}) = 5 \times \num{2.86e-4} =
\qty{1.43e-3}{mol}$.

\textbf{10.} $\num{1.43e-3} \times 55.8 = \qty{0.0798}{g} =
\qty{79.8}{mg}$.

\textbf{11.} $(80 - 79.8)/80 \approx \qty{0.3}{\percent}$.

\textbf{12.} $\qty{0.1}{mL}$ 占 \qty{14.3}{mL} 的
\qty{0.7}{\percent}，约 \qty{0.6}{mg}：结果 $79.8 \pm 0.6$ mg 与
标签一致。

\textbf{13.} $14.1$: $5 \times 0.0200 \times 0.0141 \times 55.8 =
\qty{78.7}{mg}$; $14.4$: \qty{80.4}{mg}.

\textbf{14.} $(79.8 + 78.7 + 80.4)/3 = \qty{79.6}{mg}$.

\textbf{15.} 药片本身的铁含量略有不同，
相差约 \qty{1}{\percent}。

\textbf{16.} 不宜：$V_E$ 约为 \qty{2.9}{mL}，\omterm{def:g11:titration:titration}{滴定}管
误差达 \qty{3.5}{\percent}，是之前的五倍。

\textbf{17.} 维生素 C 也会被高锰酸根氧化：
反应将不再唯一，$V_E$ 偏大，
从而使铁的含量测定值偏高。

\textbf{18.} $\num{1.43e-3} \times 151.9 = \qty{0.217}{g}$，约
\qty{217}{mg}.

\textbf{19.} $\num{1.43e-3} \times \num{6.02e23} = \num{8.6e20}$ 个\omterm{def:g9:ions:ion}{离子}。

\textbf{20.} 每片约含铁 \qty{80}{mg}，
标签标注属实。
\end{solution}
